% GENERATED FILE. Edit canonical lessons, then run npm run generate:books.
% canonical-source-hash: fnv1a64:81d699c877c068d5
% canonical-lessons: GE-C421-faehre, GE-C421-ueberfahrt, GE-C421-kabine, GE-C421-schwimmweste, GE-C421-seekrank, GE-R421-review-414-421

\chapter{Mit der Fähre}
\label{ch:ge-mit-der-fahre}
\hlchaptermodality{421}

\begin{chapteropening}
\textbf{By the end of this chapter:} \emph{I can take a ferry: the ferry, the crossing, the cabin, the life jacket and feeling seasick.}

Complete the last lesson of chapter 421: I can take a ferry: the ferry, the crossing, the cabin, the life jacket and feeling seasick.
\end{chapteropening}

\section[die Fähre]{die Fähre — the ferry}

\label{lesson:GE-C421-faehre}



\begin{quote}
Before the new one: say the German for to refund, then the German for the refund.
\end{quote}

\subsection*{You'll want to know: die Fähre}
\textbf{die Fähre} — \textquotedblleft{}the ferry\textquotedblright{}.

\begin{quote}
Die Fähre fährt jede Stunde zur Insel.
\end{quote}

\begin{quote}
\emph{The ferry goes to the island every hour.}
\end{quote}

Plural: \emph{die Fähren}.

\begin{cousinweb}
\emph{Fähre} is built on the inherited verb \emph{fahren} 'to travel' (EtymDB lists \emph{Fähre}); English \emph{ferry} comes from the same Germanic family.
\end{cousinweb}

\begin{grammarlens}[title={What you've built}]
One more word for telling, explaining and describing at B1.
\end{grammarlens}

\subsection*{Your turn}
\begin{itemize}
\raggedright
  \item \emph{Say it:} \emph{die Fähre}
  \item \emph{Say it:} \emph{die Fähre}, once more
  \item \emph{Say it:} a sentence of your own with \emph{die Fähre}
  \item \emph{Recall:} say the German for to refund, then the German for the refund, then say \emph{die Fähre} again
\end{itemize}

\begin{tcolorbox}[breakable,colback=teal!4,colframe=teal!35!black,title={Before you move on}]
What does die Fähre mean? (\textquotedblleft{}The ferry\textquotedblright{}.) Say it once more.
\end{tcolorbox}
\section[die Überfahrt]{die Überfahrt — the crossing (by sea)}

\label{lesson:GE-C421-ueberfahrt}



\begin{quote}
Before the new one: say the German for the refund, then the German for the ferry.
\end{quote}

\subsection*{You'll want to know: die Überfahrt}
\textbf{die Überfahrt} — \textquotedblleft{}the crossing (by sea)\textquotedblright{}.

\begin{quote}
Die Überfahrt dauert eine Stunde.
\end{quote}

\begin{quote}
\emph{The crossing takes an hour.}
\end{quote}

Plural: \emph{die Überfahrten}.

\begin{cousinweb}
\emph{Überfahrt} is built from \emph{über} 'over, across' and \emph{Fahrt} 'journey', from the inherited verb \emph{fahren} (EtymDB).
\end{cousinweb}

\begin{grammarlens}[title={What you've built}]
One more word for telling, explaining and describing at B1.
\end{grammarlens}

\subsection*{Your turn}
\begin{itemize}
\raggedright
  \item \emph{Say it:} \emph{die Überfahrt}
  \item \emph{Say it:} \emph{die Überfahrt}, once more
  \item \emph{Say it:} a sentence of your own with \emph{die Überfahrt}
  \item \emph{Recall:} say the German for the refund, then the German for the ferry, then say \emph{die Überfahrt} again
\end{itemize}

\begin{tcolorbox}[breakable,colback=teal!4,colframe=teal!35!black,title={Before you move on}]
What does die Überfahrt mean? (\textquotedblleft{}The crossing (by sea)\textquotedblright{}.) Say it once more.
\end{tcolorbox}
\section[die Kabine]{die Kabine — the cabin; the changing cubicle}

\label{lesson:GE-C421-kabine}



\begin{quote}
Before the new one: say the German for the ferry, then the German for the crossing.
\end{quote}

\subsection*{You'll want to know: die Kabine}
\textbf{die Kabine} — \textquotedblleft{}the cabin; the changing cubicle\textquotedblright{}.

\begin{quote}
Unsere Kabine hat ein kleines rundes Fenster.
\end{quote}

\begin{quote}
\emph{Our cabin has a small round window.}
\end{quote}

Plural: \emph{die Kabinen}. In a clothes shop, \emph{die Kabine} is the changing cubicle.

\begin{cousinweb}
\emph{Kabine} was borrowed from French \emph{cabine}, a word linked with English \emph{cabin}; both are usually traced to Late Latin \emph{capanna} 'hut'. EtymDB lists the German word without a source, so the line is hedged.
\end{cousinweb}

\begin{grammarlens}[title={What you've built}]
One more word for telling, explaining and describing at B1.
\end{grammarlens}

\subsection*{Your turn}
\begin{itemize}
\raggedright
  \item \emph{Say it:} \emph{die Kabine}
  \item \emph{Say it:} \emph{die Kabine}, once more
  \item \emph{Say it:} a sentence of your own with \emph{die Kabine}
  \item \emph{Recall:} say the German for the ferry, then the German for the crossing, then say \emph{die Kabine} again
\end{itemize}

\begin{tcolorbox}[breakable,colback=teal!4,colframe=teal!35!black,title={Before you move on}]
What does die Kabine mean? (\textquotedblleft{}The cabin; the changing cubicle\textquotedblright{}.) Say it once more.
\end{tcolorbox}
\section[die Schwimmweste]{die Schwimmweste — the life jacket}

\label{lesson:GE-C421-schwimmweste}



\begin{quote}
Before the new one: say the German for the crossing, then the German for the cabin.
\end{quote}

\subsection*{You'll want to know: die Schwimmweste}
\textbf{die Schwimmweste} — \textquotedblleft{}the life jacket\textquotedblright{}.

\begin{quote}
Unter jedem Sitz liegt eine Schwimmweste.
\end{quote}

\begin{quote}
\emph{There is a life jacket under every seat.}
\end{quote}

Plural: \emph{die Schwimmwesten}.

\begin{cousinweb}
\emph{Schwimmweste} is built from \emph{schwimmen} 'to swim, to float', inherited from Proto-Germanic \emph{swimmaną}, and \emph{Weste} 'waistcoat', which came through French and Italian \emph{veste} from Latin \emph{vestis} 'clothing' (EtymDB).
\end{cousinweb}

\begin{grammarlens}[title={What you've built}]
One more word for telling, explaining and describing at B1.
\end{grammarlens}

\subsection*{Your turn}
\begin{itemize}
\raggedright
  \item \emph{Say it:} \emph{die Schwimmweste}
  \item \emph{Say it:} \emph{die Schwimmweste}, once more
  \item \emph{Say it:} a sentence of your own with \emph{die Schwimmweste}
  \item \emph{Recall:} say the German for the crossing, then the German for the cabin, then say \emph{die Schwimmweste} again
\end{itemize}

\begin{tcolorbox}[breakable,colback=teal!4,colframe=teal!35!black,title={Before you move on}]
What does die Schwimmweste mean? (\textquotedblleft{}The life jacket\textquotedblright{}.) Say it once more.
\end{tcolorbox}
\section[seekrank]{seekrank — seasick}

\label{lesson:GE-C421-seekrank}



\begin{quote}
Before the new one: say the German for the cabin, then the German for the life jacket.
\end{quote}

\subsection*{You'll want to know: seekrank}
\textbf{seekrank} — \textquotedblleft{}seasick\textquotedblright{}.

\begin{quote}
Bei hohen Wellen wird sie schnell seekrank.
\end{quote}

\begin{quote}
\emph{When the waves are high she quickly gets seasick.}
\end{quote}

\emph{Die Seekrankheit} — seasickness.

\begin{cousinweb}
\emph{seekrank} is built from \emph{See} 'sea, lake' and \emph{krank} 'ill', both inherited: \emph{krank} comes through Middle High German from Old High German \emph{krank} (EtymDB), a word usually said to have first meant 'weak'.
\end{cousinweb}

\begin{grammarlens}[title={What you've built}]
One more word for telling, explaining and describing at B1.
\end{grammarlens}

\subsection*{Your turn}
\begin{itemize}
\raggedright
  \item \emph{Say it:} \emph{seekrank}
  \item \emph{Say it:} \emph{seekrank}, once more
  \item \emph{Say it:} a sentence of your own with \emph{seekrank}
  \item \emph{Recall:} say the German for the cabin, then the German for the life jacket, then say \emph{seekrank} again
\end{itemize}

\begin{tcolorbox}[breakable,colback=teal!4,colframe=teal!35!black,title={Before you move on}]
What does seekrank mean? (\textquotedblleft{}Seasick\textquotedblright{}.) Say it once more.
\end{tcolorbox}
\section[(dialogue)]{Words from chapters 414–421 — a review}

\label{lesson:GE-R421-review-414-421}



\begin{quote}
Say the German for the connection and the local transport.
\end{quote}

\subsection*{Your turn}
Say these aloud:

\begin{itemize}
\raggedright
  \item the connection, the local transport, the long-distance transport, to commute, the commuter
  \item the locker, the left-luggage office, the escalator, to ride without a ticket, to queue
  \item to check in, the boarding pass, the gate, the hand luggage, to give up
  \item the security check, to search, the customs, to declare, the passport
  \item the flight attendant, the window seat, the aisle seat, to take off, the landing
  \item the stopover, the time difference, to be late, to be stuck, to make an emergency landing
  \item to cancel, the cancellation, to rebook, to refund, the refund
  \item the ferry, the crossing, the cabin, the life jacket, seasick
\end{itemize}

\begin{tcolorbox}[breakable,colback=teal!4,colframe=teal!35!black,title={Before you move on}]
The connection? (\textbf{der Anschluss}.) The local transport? (\textbf{der Nahverkehr}.) The long-distance transport? (\textbf{der Fernverkehr}.) To commute? (\textbf{pendeln}.) The commuter? (\textbf{der Pendler}.) The locker? (\textbf{das Schließfach}.) The left-luggage office? (\textbf{die Gepäckaufbewahrung}.) The escalator? (\textbf{die Rolltreppe}.) To ride without a ticket? (\textbf{schwarzfahren}.) To queue? (\textbf{anstellen}.) To check in? (\textbf{einchecken}.) The boarding pass? (\textbf{die Bordkarte}.) The gate? (\textbf{der Flugsteig}.) The hand luggage? (\textbf{das Handgepäck}.) To give up? (\textbf{aufgeben}.) The security check? (\textbf{die Sicherheitskontrolle}.) To search? (\textbf{durchsuchen}.) The customs? (\textbf{der Zoll}.) To declare? (\textbf{verzollen}.) The passport? (\textbf{der Reisepass}.) The flight attendant? (\textbf{der Flugbegleiter}.) The window seat? (\textbf{der Fensterplatz}.) The aisle seat? (\textbf{der Gangplatz}.) To take off? (\textbf{abfliegen}.) The landing? (\textbf{die Landung}.) The stopover? (\textbf{die Zwischenlandung}.) The time difference? (\textbf{die Zeitverschiebung}.) To be late? (\textbf{verspäten}.) To be stuck? (\textbf{festsitzen}.) To make an emergency landing? (\textbf{notlanden}.) To cancel? (\textbf{stornieren}.) The cancellation? (\textbf{die Stornierung}.) To rebook? (\textbf{umbuchen}.) To refund? (\textbf{erstatten}.) The refund? (\textbf{die Erstattung}.) The ferry? (\textbf{die Fähre}.) The crossing? (\textbf{die Überfahrt}.) The cabin? (\textbf{die Kabine}.) The life jacket? (\textbf{die Schwimmweste}.) Seasick? (\textbf{seekrank}.)
\end{tcolorbox}
